\section{Collection platform}
\label{sec:platform}

The platform is a dual five-degree-of-freedom rig with two teleoperation
interfaces and a multi-view sensor suite. This section walks the arms, the
operator interfaces, the sensing, and the host and buses, and it ends with
the fixed state the rig must reach before collection begins.

\paragraph{Arms.} Two follower arms, each with five actuated body joints
and a gripper, sit in a fixed left and right layout. That layout is shared
with the digital twin used for earlier simulation work, so the physical
placement and the simulated placement come from the same description. Each
arm carries a camera at the wrist that looks along the gripper.

\paragraph{Operator interfaces.} The rig accepts two kinds of
teleoperation, and both produce demonstrations with identical recorded
semantics. A head-mounted controller drives the arms in Cartesian space
through an inverse-kinematics solver. A pair of leader arms, kinematically
identical to the followers, drives them directly in joint space. The choice
of interface changes how the operator moves, not what the recorder stores.

\paragraph{Sensing.} Each gripper carries two visual-tactile fingertip
cameras in addition to the wrist camera, and a colour camera overlooks the
workspace from above; a depth-capable camera may take that overhead position
instead, adding an aligned-depth image. The colour images, together with
that optional overhead depth, are the exterior views a policy sees; the
fingertip views add local contact information. All views are recorded
through one pipeline so that no stream is privileged over another.

\paragraph{Host and buses.} Every device reaches a single laptop. The
cameras and arms share two powered universal-serial-bus hubs, one per side,
and the head-mounted controller uses a direct link. Shared-bus bandwidth is
a genuine constraint when several compressed image streams run at once. We
do not assume the bus copes; instead the recorder measures each stream's
rate and dropped frames, so a saturated bus is visible rather than silent
(Section~\ref{sec:collection}).

\paragraph{Final-state readiness.} Collection is only meaningful when the
rig starts from a fixed, known state: every arm calibrated and bound to a
stable port alias, the wrist joints limited so the fingertip cables cannot
be twisted, the arms at their defined home pose, and every camera opened at
its configured resolution. Each colour camera is likewise bound to a stable
per-socket alias rather than an enumeration index, the overhead colour
camera included (a depth-capable overhead camera is pinned by its own serial
instead), so replugging a hub can never silently swap one view for another
between collection and deployment; the same bound identities drive both the
live monitor and the recorder, and the monitor shows every stream before a
session begins, including the overhead view, with its aligned depth rendered
alongside when a depth camera is fitted. The same monitor also
replays a recorded episode from the stored frames, so a session can be
reviewed exactly as it was seen during collection, and that review can be
exported as a video. A
preflight check confirms these conditions and
the host-tuning items that reduce timing jitter, and reports a single
ready-or-not verdict. The platform's fixed dimensions and camera
calibration are measured and stored so the rig can be rebuilt
(Section~\ref{sec:replicability}).
